\reading{MR-8}{Correlation Basics: Definitions, Applications, and Terminology}{strike}{7}
% ==========================================================

\begin{mrkeybox}[title={Learning objectives in this reading}]
\small
\begin{enumerate}[label=\textbf{\alph*.},leftmargin=1.7em]
\item Describe financial correlation risk and the areas in which it appears in finance.
\item Explain how correlation contributed to the global financial crisis of 2007--2009.
\item Describe how correlation impacts the price of quanto options as well as other multi-asset
      exotic options.
\item Describe the structure, uses, and payoffs of a correlation swap.
\item Estimate the impact of different correlations between assets in the trading book on the
      VaR capital charge.
\item Explain the role of correlation risk in market risk and credit risk.
\item Explain how correlation risk relates to systemic and concentration risk.
\end{enumerate}
\end{mrkeybox}

% ---------------------------------------------------------
\lo{MR-8 a}{Describe financial correlation risk and the areas in which it appears in finance.}

\subsection{Correlation basics}
Correlation explains the degree of relationship between two variables, and it is measured
through correlation analysis. The measure of correlation is called the \term{correlation
coefficient}. The degree of relationship is expressed by correlation, which ranges from $-1$ to
$+1$.

\begin{mrdefbox}[title={Financial correlation risk}]
Financial correlation risk is the risk of financial loss due to adverse movements in correlation
between two or more variables --- the risk that the relationship between different financial or
non-financial assets changes unexpectedly, leading to potential financial losses. These
variables can comprise any financial variables. As the correlation of asset returns increases,
the risk of financial loss, which is often measured by the VaR concept, also increases.
\end{mrdefbox}

\subsection{Where it shows up}

\begin{center}
\small
\begin{tabularx}{\linewidth}{@{}p{3.5cm} >{\raggedright\arraybackslash}X@{}}
\arrayrulecolor{FRMRule}\toprule
\rowcolor{FRMNavy} \hdr{Financial instruments} & \hdr{}\\
Interest rates and commodity prices & Typically there is a \textit{negative} correlation. If
  interest rates rise, commodity prices might fall, hurting commodity investors.\\
\rowcolor{FRMSoft}
Greek and Mexican bonds & During the 2012 Greek crisis, bonds in seemingly unrelated markets
  also suffered, through a sudden \textit{positive} correlation with Greek bonds.\\
The 2007 crisis & Showed how globally diversified portfolios could suffer when diverse assets
  suddenly correlate positively in a downturn.\\
\midrule
\rowcolor{FRMNavy} \hdr{Non-financial} & \hdr{}\\
Currency devaluation & In 2012, US exporters experienced losses due to devaluation of the euro
  during the sovereign debt crisis.\\
\rowcolor{FRMSoft}
Growth linkages & Low GDP growth for the US had major adverse impacts on Asian and European
  exporters.\\
Political events & Uprisings in the Middle East can impact airline travel through rising oil
  prices, affecting industries and investors in the travel sector.\\
\bottomrule
\end{tabularx}
\end{center}

\subsection{Two categories of financial correlation}

\begin{center}
\small
\begin{tabularx}{\linewidth}{@{}p{3.2cm} >{\raggedright\arraybackslash}X@{}}
\arrayrulecolor{FRMRule}\toprule
\rowcolor{FRMNavy} \hdr{} & \hdr{Definition and examples}\\
\textbf{1. Static}\newline correlations
  & Measure how two or more financial assets are associated \textit{at a certain point in time}
    or within a certain time period. Fixed over the period.
    \begin{itemize}[topsep=2pt]
    \item Correlating bond prices and their respective yields at a point in time.
    \item The classic VaR model --- you calculate VaR for a certain period of time.
    \item The copula approach for CDOs, measuring default correlations between all assets in
          the CDO for a certain time period.
    \item The binomial default correlation model.
    \end{itemize}\\
\rowcolor{FRMSoft}
\textbf{2. Dynamic}\newline correlations
  & Measure how two or more financial assets \textit{move together in time}. Changing over time.
    \begin{itemize}[topsep=2pt]
    \item ``Pairs trading'' --- where one asset is purchased and another is sold.
    \item Within the deterministic correlation approaches, the Heston model (1993) correlates
          the Brownian motions of assets.
    \item Correlations behave randomly and unpredictably (stochastic correlation processes).
    \end{itemize}\\
\bottomrule
\end{tabularx}
\end{center}

\begin{mrkeybox}[title={Points just to consider}]
\begin{enumerate}[label=\textbf{\alph*)}]
\item Standard deviation and variance calculations help determine the risk of a portfolio.
\item The focus is on risk-adjusted returns to assess portfolio performance.
\item The return/risk ratio measures the average return of a portfolio relative to its risk,
      influenced by correlation.
\end{enumerate}
\end{mrkeybox}

\subsection{Correlation risk seen through a CDS}

\begin{center}
\begin{tikzpicture}[every node/.style={font=\scriptsize\sffamily},
  party/.style={rectangle, rounded corners=3pt, line width=0.9pt,
    text width=3.3cm, align=center, inner sep=7pt, minimum height=1.6cm}]
  \node[party, draw=FRMNavy!65, fill=PaleNavy] (inv) at (0,0)
    {\textbf{Investor $i$}\\ the credit default swap \textit{buyer}};
  \node[party, draw=FRMGold!70, fill=FRMPaper] (cpty) at (7.4,0)
    {\textbf{Counterparty $C$}\\ the CDS \textit{seller}\\ (BNP Paribas)};
  \node[party, draw=FRMRed!60, fill=PaleRed] (ref) at (0,-3.1)
    {\textbf{Reference entity $r$}\\ the reference asset\\ (Spain)};

  \draw[FRMNavy, -{Stealth[length=5pt]}, line width=0.9pt]
    ([yshift=7pt]inv.east) -- node[above, text=black]{Fixed CDS spread $s$}
    ([yshift=7pt]cpty.west);
  \draw[FRMGold, dashed, -{Stealth[length=5pt]}, line width=0.9pt]
    ([yshift=-9pt]cpty.west) -- node[below, text=black, align=center]
    {USD 1m payout\\ if $r$ defaults} ([yshift=-9pt]inv.east);
  \draw[FRMRed, -{Stealth[length=5pt]}, line width=0.9pt]
    ([xshift=-14pt]inv.south) -- node[left, text=black]{USD 1m} ([xshift=-14pt]ref.north);
  \draw[FRMRed, -{Stealth[length=5pt]}, line width=0.9pt]
    ([xshift=14pt]ref.north) -- node[right, text=black]{coupon $k$} ([xshift=14pt]inv.south);

  \draw[FRMGreen!70, dotted, line width=1.3pt] (ref.east) -- (cpty.south west);
  \node[anchor=west, text=FRMGreen, align=left, fill=white, inner sep=2pt]
    at (5.35,-2.35) {the correlation that\\ actually matters: $\rho(r, C)$};
\end{tikzpicture}
\end{center}
\figcap{An investor hedging Spanish bond exposure with a CDS. The dotted line is the leg the
notes are really about: the hedge only works if Spain and BNP Paribas do not go down together.}

\begin{mrdefbox}[title={What the fixed CDS spread depends on}]
\begin{enumerate}
\item The reference asset's default probability.
\item The joint default correlation between $r$ and BNP Paribas.
\end{enumerate}
\medskip
If correlation \term{increases}, then the value of the CDS (or spread) will \term{decrease},
because if both Spain and BNP Paribas default you get nothing --- this is \term{wrong-way risk}.
\end{mrdefbox}

\begin{mrtrapbox}[title={Non-monotonic}]
The relationship between the CDS spread and correlation may be \term{non-monotonous}. If
correlation increases from $+1$ to $-0.4$, the CDS spread may decrease slightly.
\end{mrtrapbox}

\subsection{Correlations in financial investments}
We know from the Sharpe ratio that increased diversification increases risk-return performance:
high diversification means low correlation between assets. In 1952, Harry Markowitz provided the
foundation of modern investment theory by demonstrating the role that correlation plays in
reducing risk.

\begin{mrfmlbox}
\[ \mu_P \;=\; w_X \mu_X + w_Y \mu_Y \qquad\qquad
   \sigma_P \;=\; \sqrt{ w_X^2 \sigma_X^2 + w_Y^2 \sigma_Y^2 + 2\,w_X w_Y \,\mathrm{cov}_{XY} } \]
\mrvk{$w_X$, $w_Y$ = portfolio weights;\quad $\mu_X$, $\mu_Y$ = mean returns;\quad
$\sigma_X$, $\sigma_Y$ = standard deviations;\quad $\mathrm{cov}_{XY}$ = covariance between the
two assets. Only the covariance term carries the correlation, and only that term can be
negative.}{}
\end{mrfmlbox}

\begin{mrfmlbox}
\[ \mathrm{cov}_{XY} \;=\; \frac{\sum_{t=1}^{n} (X_t - \mu_X)(Y_t - \mu_Y)}{n-1}
   \qquad\qquad
   \rho_{XY} \;=\; \frac{\mathrm{cov}_{XY}}{\sigma_X \sigma_Y} \]
\mrvk{$n$ = number of return observations. Dividing the covariance by the two standard deviations
is what rescales it onto the $[-1, +1]$ range.}{}
\end{mrfmlbox}

\begin{center}
\small
\begin{tabular}{c r r r r}
\arrayrulecolor{FRMRule}\toprule
\rowcolor{FRMNavy} \hdr{Year} & \hdr{$X$} & \hdr{$Y$} & \hdr{Return $X$} & \hdr{Return $Y$}\\
2009 & 90 & 150 & & \\
\rowcolor{FRMSoft} 2010 & 120 & 180 & 0.3333 & 0.2000\\
2011 & 105 & 340 & $(0.1250)$ & 0.8889\\
\rowcolor{FRMSoft} 2012 & 170 & 320 & 0.6190 & $(0.0588)$\\
2013 & 150 & 360 & $(0.1176)$ & 0.1250\\
\rowcolor{FRMSoft} 2014 & 270 & 310 & 0.8000 & $(0.1389)$\\
\midrule
\multicolumn{3}{r}{\textit{Average return}} & \textbf{0.3019} & \textbf{0.2032}\\
\bottomrule
\end{tabular}
\end{center}

\begin{center}
\scriptsize
\begin{tabular}{c r r r r r r r}
\arrayrulecolor{FRMRule}\toprule
\rowcolor{FRMNavy} \hdr{Year} & \hdr{Ret.\ $X$} & \hdr{Ret.\ $Y$} & \hdr{$X_t-\mu_X$}
  & \hdr{$Y_t-\mu_Y$} & \hdr{$(X_t-\mu_X)^2$} & \hdr{$(Y_t-\mu_Y)^2$}
  & \hdr{cross product}\\
2010 & 0.3333 & 0.2000 & 0.0314 & $(0.0032)$ & 0.0010 & 0.0000 & $(0.0001)$\\
\rowcolor{FRMSoft} 2011 & $(0.1250)$ & 0.8889 & $(0.4269)$ & 0.6857 & 0.1823 & 0.4701 & $(0.2927)$\\
2012 & 0.6190 & $(0.0588)$ & 0.3171 & $(0.2621)$ & 0.1006 & 0.0687 & $(0.0831)$\\
\rowcolor{FRMSoft} 2013 & $(0.1176)$ & 0.1250 & $(0.4196)$ & $(0.0782)$ & 0.1761 & 0.0061 & 0.0328\\
2014 & 0.8000 & $(0.1389)$ & 0.4981 & $(0.3421)$ & 0.2481 & 0.1170 & $(0.1704)$\\
\midrule
\textit{Mean / sum} & \textbf{0.3019} & \textbf{0.2032} & & & 0.7079 & 0.6620 & $(0.5135)$\\
\multicolumn{5}{r}{\textit{Variance}} & \textbf{0.1770} & \textbf{0.1655} & $(0.1284)$\\
\multicolumn{5}{r}{\textit{Standard deviation}} & \textbf{0.4207} & \textbf{0.4068} & \\
\multicolumn{5}{r}{\textit{Correlation}} & \multicolumn{2}{c}{$\mathbf{(0.7501)}$}\\
\bottomrule
\end{tabular}
\end{center}

\[ \rho_{XY} \;=\; \frac{-0.1284}{0.4207 \times 0.4068} \;=\; \mathbf{-0.7501} \]

\begin{center}
\begin{tikzpicture}
\begin{axis}[
  width=11.0cm, height=5.8cm,
  xmin=-1, xmax=1, ymin=0, ymax=3.0,
  xlabel={\scriptsize Correlation $\rho_{XY}$}, ylabel={\scriptsize Return / Risk},
  xlabel style={font=\scriptsize}, ylabel style={font=\scriptsize},
  tick label style={font=\scriptsize},
  xtick={-0.9,-0.7,-0.5,-0.3,-0.1,0.1,0.3,0.5,0.7,0.9},
  ytick={0,0.5,1.0,1.5,2.0,2.5,3.0},
  axis line style={draw=black!55}, clip=false]
\addplot[FRMNavy, line width=1.2pt, smooth] coordinates {
(-0.9,2.7225)(-0.8,1.9278)(-0.7,1.5748)(-0.6,1.3641)(-0.5,1.2203)(-0.4,1.1141)
(-0.3,1.0315)(-0.2,0.9649)(-0.1,0.9098)(0.0,0.8631)(0.1,0.8230)(0.2,0.7879)
(0.3,0.7570)(0.4,0.7295)(0.5,0.7048)(0.6,0.6824)(0.7,0.6620)(0.8,0.6434)
(0.9,0.6262)(1.0,0.6104)};
\draw[FRMRed, dashed, line width=0.7pt] (axis cs:-0.7501,0) -- (axis cs:-0.7501,1.72);
\node[anchor=west, font=\scriptsize\sffamily, text=FRMRed, fill=white, inner sep=1.5pt] at (axis cs:-0.70,1.85)
  {this portfolio, $\rho = -0.75$};
\end{axis}
\end{tikzpicture}
\end{center}
\figcap{Built from the $X$ and $Y$ statistics in the table above at equal weights. The ratio
barely moves across the whole positive half of the range, then climbs steeply once correlation
turns strongly negative. Diversification pays almost nothing until correlation goes negative,
and then it pays enormously.}

% ---------------------------------------------------------
\lo{MR-8 b}{Explain how correlation contributed to the global financial crisis of 2007--2009.}

\begin{enumerate}
\item \term{Misinformed risk models.} Risk managers and investors underestimated the
      correlations between different financial assets, especially in complex structured products
      like collateralized debt obligations (CDOs). This led to a false sense of security and
      mispricing of risk.
\item \term{CDO tranche mismanagement.} Large hedge funds believed that they could hedge risk by
      being \textit{long} the riskiest CDO tranches (equity) and \textit{short} the less risky
      tranches (mezzanine). However, they did not fully grasp the correlation dynamics across
      these tranches, resulting in substantial losses when the crisis hit.
\item \term{Bond market impact.} The downgrading of bonds issued by U.S.\ automobile makers to
      junk status, and the subsequent correlation shift among different bond tranches,
      contributed to market turmoil. Correlations between bonds increased, causing price
      declines and impacting investors.
\item \term{CDO market growth.} The rapid growth of the CDO market, linked to overvalued real
      estate and subprime mortgages, made it susceptible to severe losses when the housing
      bubble burst. High correlations between these assets exacerbated the crisis.
\item \term{Credit default swap market.} The CDS market also grew significantly, but some
      institutions like AIG were not adequately prepared to cover their CDS liabilities, leading
      to financial instability.
\end{enumerate}

\begin{mrkeybox}[title={Remember}]
The underestimation of correlation risk, particularly in complex financial products, and the
interconnectedness of various markets amplified the impact of the financial crisis.
\end{mrkeybox}

% ---------------------------------------------------------
\lo{MR-8 c}{Describe how correlation impacts the price of quanto options as well as other multi-asset exotic options.}

Correlation trading strategies involve trading assets that have prices determined by the
comovement of one or more assets over time. \term{Correlation options} have prices that are very
sensitive to the correlation between two assets and are often referred to as \term{multi-asset
options}.

\begin{center}
\small
\begin{tabularx}{\linewidth}{@{}p{4.4cm} >{\raggedright\arraybackslash}X@{}}
\arrayrulecolor{FRMRule}\toprule
\rowcolor{FRMNavy} \hdr{Strategy} & \hdr{Payoff}\\
a) Option on the higher of two stocks & Pays the maximum of the two stock prices,
  $\max(S_1, S_2)$.\\
\rowcolor{FRMSoft}
b) Exchange option & Pays the difference between the prices of two assets.\\
c) Spread call option & Pays the difference between two stock prices minus the strike price.\\
\bottomrule
\end{tabularx}
\end{center}

\begin{mrtrapbox}[title={Remember the direction}]
\term{Lower correlation often leads to higher option prices.} A lower correlation means one
asset might go up while the other goes down, increasing the chance of a profitable outcome.
\end{mrtrapbox}

\subsection{Quanto options}
A quanto option allows a domestic investor to exchange his potential option payoff in a foreign
currency back into his home currency \term{at a fixed exchange rate}. Low correlation means a
high quanto price.

\begin{mrexbox}[title={Example --- a quanto call on the Nikkei}]
An American believes the Nikkei will increase, but she is worried about a decreasing yen. The
investor can buy a quanto call on the Nikkei, with the yen payoff being converted into dollars
at a fixed exchange rate.
\end{mrexbox}

% ---------------------------------------------------------
\lo{MR-8 d}{Describe the structure, uses, and payoffs of a correlation swap.}

A correlation swap enables trading a \term{fixed} correlation rate between two or more assets
against the \term{actual or realized} correlation.

\begin{center}
\begin{tikzpicture}[every node/.style={font=\scriptsize\sffamily},
  side/.style={rectangle, rounded corners=3pt, line width=0.9pt,
    text width=3.6cm, align=center, inner sep=7pt, minimum height=1.5cm}]
  \node[side, draw=FRMGreen!65, fill=PaleGreen] (buy) at (0,0)
    {\textit{Buying correlation}\\ \textbf{fixed rate $\rho$ payer}};
  \node[side, draw=FRMRed!60, fill=PaleRed] (sell) at (7.4,0)
    {\textit{Selling correlation}\\ \textbf{fixed rate $\rho$ receiver}};
  \draw[FRMNavy, -{Stealth[length=5pt]}, line width=0.9pt]
    ([yshift=8pt]buy.east) -- node[above, text=black]{fixed $\rho = 30\%$}
    ([yshift=8pt]sell.west);
  \draw[FRMGold, -{Stealth[length=5pt]}, line width=0.9pt]
    ([yshift=-8pt]sell.west) -- node[below, text=black]{realized $\rho$}
    ([yshift=-8pt]buy.east);
  \node[anchor=north, text=FRMGreen, align=center] at (buy.south)
    {\\ expects correlation \textbf{to increase}};
  \node[anchor=north, text=FRMRed, align=center] at (sell.south)
    {\\ expects correlation \textbf{to decrease}};
\end{tikzpicture}
\end{center}
\figcap{The buyer pays a fixed correlation and receives whatever correlation actually turns out
to be, so the buyer profits when correlation rises. It is structurally identical to any other
fixed-for-floating swap, with correlation in place of an interest rate.}

\subsection{To calculate the payoff (for the buyer)}
\begin{enumerate}
\item Figure out what the realized (stochastic) correlation is.
\item Calculate the payoff.
\end{enumerate}

\begin{mrfmlbox}
\[ \rho_{\text{realized}} \;=\; \frac{2}{n^2 - n} \sum_{i>j} \rho_{i,j}
   \qquad\qquad
   \text{Payoff} \;=\; \text{notional} \times
   \left( \rho_{\text{realized}} - \rho_{\text{fixed}} \right) \]
\mrvk{$n$ = number of assets in the portfolio;\quad $\rho_{i,j}$ = the realized pairwise
correlation between assets $i$ and $j$, counted once per pair ($i > j$);\quad
$(n^2-n)/2$ = the number of distinct pairs, which is why the factor is $2/(n^2-n)$.}{}
\end{mrfmlbox}

\begin{mrexbox}[title={Example --- correlation swap payoff}]
A correlation swap buyer pays a fixed correlation rate of 0.2 with a notional value of
\$1 million for one year for a portfolio of \textbf{three} assets. The realized pairwise
correlations of the daily log returns $[\ln(S_t/S_{t-1})]$ at maturity for the three assets are
$\rho_{2,1} = 0.6$, $\rho_{3,1} = 0.2$, and $\rho_{3,2} = 0.04$. (Note that for all pairs
$i > j$.) What is the correlation swap buyer's payoff?
\tcblower
\small
The realized correlation is calculated as:
\[ \rho_{\text{realized}} \;=\; \frac{2}{3^2 - 3} \times (0.6 + 0.2 + 0.04) \;=\; 0.28 \]
The payoff for the correlation swap buyer is then:
\[ \$1{,}000{,}000 \times (0.28 - 0.20) \;=\; \mathbf{\$80{,}000} \]
\end{mrexbox}

% ---------------------------------------------------------
\lo{MR-8 e}{Estimate the impact of different correlations between assets in the trading book on the VaR capital charge.}

\noindent\gapbadge\hspace{4pt}{\sffamily\fontsize{7.6}{9}\selectfont\color{black!55}%
The source notes redirect this objective to the Investment Management book and develop nothing
here. It is written below from the GARP curriculum so the objective is not left blank.}\par
\vspace{6pt}

\begin{mrgapbox}[title={What the objective asks for}]
Given a trading-book portfolio, compute VaR under one correlation assumption, then re-compute it
under a different correlation and state the direction and size of the change. The correlation
enters VaR only through the covariance matrix.
\end{mrgapbox}

\subsection{Portfolio VaR with a covariance matrix}

\begin{mrfmlbox}
\[ \mathrm{VaR}_P \;=\; \sigma_P \times \alpha \times \sqrt{x}
   \qquad\qquad
   \sigma_P \;=\; \sqrt{ \beta_h \, C \, \beta_v } \]
\mrvk{$\sigma_P$ = volatility of the portfolio, which is where the correlation between the assets
enters;\quad $\alpha$ = the standard normal deviate at the chosen confidence level, so 2.3264 at
99\%;\quad $x$ = the time horizon in days;\quad $\beta_h$ = the horizontal vector of invested
amounts (price $\times$ quantity);\quad $\beta_v$ = the same vector written vertically;\quad
$C$ = the covariance matrix of the asset returns.}{}
\end{mrfmlbox}

Each entry of the covariance matrix is built from the pairwise correlation and the two
volatilities:
\[ \mathrm{Cov}_{ij} \;=\; \rho_{ij}\, \sigma_i\, \sigma_j \]

\begin{mrexbox}[title={Example --- 10-day VaR for a two-asset trading book}]
What is the 10-day VaR for a two-asset portfolio with a correlation coefficient of 0.7, a daily
standard deviation of returns of 2\% for asset 1 and 1\% for asset 2, with USD 10 million
invested in asset 1 and USD 5 million in asset 2, at a 99\% confidence level?
\tcblower
\small
\textbf{Step 1 --- build the covariance matrix.}
\begin{align*}
\mathrm{Cov}_{11} &= 1 \times 0.02 \times 0.02 = 0.0004 &
\mathrm{Cov}_{12} &= 0.7 \times 0.02 \times 0.01 = 0.00014\\
\mathrm{Cov}_{21} &= 0.7 \times 0.01 \times 0.02 = 0.00014 &
\mathrm{Cov}_{22} &= 1 \times 0.01 \times 0.01 = 0.0001
\end{align*}
\[ C \;=\; \begin{pmatrix} 0.0004 & 0.00014 \\ 0.00014 & 0.0001 \end{pmatrix} \]

\textbf{Step 2 --- multiply out $\beta_h C$.}
\[ (10 \;\; 5)\begin{pmatrix} 0.0004 & 0.00014 \\ 0.00014 & 0.0001 \end{pmatrix}
   = (0.0047 \;\; 0.0019) \]

\textbf{Step 3 --- multiply by $\beta_v$ to get the portfolio variance.}
\[ (0.0047 \;\; 0.0019)\begin{pmatrix} 10 \\ 5 \end{pmatrix}
   = 10 \times 0.0047 + 5 \times 0.0019 = 0.0565 = 5.65\% \]

\textbf{Step 4 --- take the root for $\sigma_P$.}
\[ \sigma_P = \sqrt{0.0565} = 23.77\% \]

\textbf{Step 5 --- scale by the deviate and the horizon.}
\[ \mathrm{VaR}_P = 0.2377 \times 2.3264 \times \sqrt{10} = \mathbf{USD\ 1.7486 \text{ million}} \]

\medskip
\textbf{Interpretation.} We are 99\% certain that we will not lose more than USD 1.7486 million
in the next 10 days due to correlated market price changes of assets 1 and 2.
\end{mrexbox}

\subsection{The impact of changing the correlation}

\begin{center}
\begin{tikzpicture}
\begin{axis}[
  width=11.0cm, height=5.4cm,
  xmin=-1.05, xmax=1.05, ymin=1.0, ymax=1.95,
  xlabel={\scriptsize Correlation between asset 1 and asset 2},
  ylabel={\scriptsize 10-day 99\% VaR (USD m)},
  xlabel style={font=\scriptsize}, ylabel style={font=\scriptsize},
  tick label style={font=\scriptsize},
  xtick={-1,-0.75,-0.5,-0.25,0,0.25,0.5,0.75,1},
  ytick={1.0,1.2,1.4,1.6,1.8},
  axis line style={draw=black!55}, clip=false]
\addplot[FRMNavy, line width=1.2pt, smooth] coordinates {
(-1.0,1.1035)(-0.9,1.1515)(-0.8,1.1976)(-0.7,1.2420)(-0.6,1.2848)(-0.5,1.3263)
(-0.4,1.3665)(-0.3,1.4055)(-0.2,1.4435)(-0.1,1.4805)(0.0,1.5166)(0.1,1.5519)
(0.2,1.5864)(0.3,1.6202)(0.4,1.6532)(0.5,1.6856)(0.6,1.7174)(0.7,1.7487)
(0.8,1.7794)(0.9,1.8095)(1.0,1.8392)};
\draw[FRMRed, dashed, line width=0.7pt] (axis cs:0.7,1.0) -- (axis cs:0.7,1.7487);
\node[circle, fill=FRMRed, inner sep=1.7pt] at (axis cs:0.7,1.7487) {};
\node[anchor=east, font=\scriptsize\sffamily, text=FRMRed, fill=white, inner sep=1.5pt] at (axis cs:0.62,1.79)
  {worked example: $\rho = 0.7 \Rightarrow 1.75$};
\node[circle, fill=FRMNavy, inner sep=1.5pt] at (axis cs:-1,1.1035) {};
\node[anchor=south west, font=\scriptsize\sffamily, text=FRMNavy] at (axis cs:-0.96,1.02)
  {$\rho = -1 \Rightarrow 1.10$};
\end{axis}
\end{tikzpicture}
\end{center}
\figcap{The same two positions, the same volatilities, the same confidence level. Only the
correlation changes. VaR runs from about USD 1.1 million at perfect negative correlation to
about USD 1.9 million at perfect positive correlation, and the curve is concave, so the risk
reduction accelerates as correlation falls.}

\begin{mrkeybox}[title={The direction to state out loud}]
The \term{lower} the correlation, the \term{lower} the risk measured by VaR, and so the lower
the capital charge. Preferably the correlation is negative: if one asset decreases, the other on
average increases, reducing the overall risk. The impact is strong, roughly a 70\% increase in
VaR across the full correlation range in this example.
\end{mrkeybox}

% ---------------------------------------------------------
\lo{MR-8 f}{Explain the role of correlation risk in market risk and credit risk.}

\subsection{Correlation risk and market risk}
\begin{enumerate}[label=\textbf{\alph*)}]
\item Market risk encompasses various factors affecting asset prices, like interest rates,
      currency fluctuations, and equity or commodity prices.
\item VaR and expected shortfall are risk measures that \textit{consider correlations} in asset
      prices.
\item Inaccurate correlation estimates in VaR or ES calculations can lead to
      \term{underestimating} market risk.
\end{enumerate}

\subsection{Correlation risk and credit risk}
\begin{enumerate}[label=\textbf{\alph*)}]
\item Credit risk deals with the potential for borrowers to default on loans.
\item \term{Migration risk} is the downgrading of a borrower's credit rating, increasing default
      risk.
\item \term{Default correlation} measures how likely multiple borrowers default simultaneously.
\item \term{Lower} default correlation allows for \textit{better} diversification of credit risk
      across lenders.
\item The financial crisis of 2007--2009 is an example of how increased correlation between
      seemingly uncorrelated assets (like mortgages and bonds) can amplify credit risk.
\end{enumerate}

% ---------------------------------------------------------
\lo{MR-8 g}{Explain how correlation risk relates to systemic and concentration risk.}

\subsection{Correlation risk and systemic risk}
\begin{enumerate}[label=\textbf{\alph*)}]
\item Systemic risk refers to the potential collapse of the entire financial system.
\item During a crisis, correlations between \textit{all} asset classes tend to increase, making
      diversification less effective.
\item The 2007 crisis exemplified this, with correlations between U.S.\ stocks, bonds, and
      international equities rising significantly.
\end{enumerate}

\subsection{Concentration risk and correlation risk}
\begin{enumerate}[label=\textbf{\alph*)}]
\item \term{Concentration risk} is the risk of excessive exposure to a single borrower or a
      specific industry.
\item A \term{concentration ratio} measures the diversification of loan portfolios.
\end{enumerate}

\begin{center}
\begin{tikzpicture}[every node/.style={font=\scriptsize\sffamily},
  chip/.style={rectangle, rounded corners=2.5pt, draw=FRMNavy!45, fill=PaleNavy,
    line width=0.7pt, text width=2.5cm, align=center, inner sep=5pt, minimum height=0.95cm}]
  \node[chip] (a) at (0,0) {Concentration\\ \textbf{ratio} $\downarrow$};
  \node[chip] (b) at (3.35,0) {Concentration\\ \textbf{risk} $\downarrow$};
  \node[chip] (c) at (6.70,0) {\textbf{Diversification}\\ $\uparrow$};
  \node[chip] (d) at (10.05,0) {\textbf{Correlation}\\ $\downarrow$};
  \foreach \x/\y in {a/b, b/c, c/d}{
    \draw[FRMGold, -{Stealth[length=4.5pt]}, line width=1pt] (\x) -- (\y);}
\end{tikzpicture}
\end{center}
\figcap{All four move together. A lower concentration ratio means the book is spread across more
borrowers, which cuts concentration risk, raises diversification, and shows up as lower
effective correlation in the portfolio.}

\begin{mrfmlbox}
\[ P(AB) \;=\; \rho_{AB} \sqrt{ PD_A\,(1 - PD_A) \times PD_B\,(1 - PD_B) }
   \;+\; PD_A \times PD_B \]
\mrvk{$P(AB)$ = the joint probability that both $A$ and $B$ default;\quad $\rho_{AB}$ = the
default correlation between them;\quad $PD_A$, $PD_B$ = the individual default probabilities.
When $\rho_{AB} = 0$ the first term vanishes and the joint probability collapses to the simple
product $PD_A \times PD_B$.}{}
\end{mrfmlbox}

\begin{mrexbox}[title={Example --- EL, concentration ratio and joint default probability}]
Commercial Bank Y has loans of \$2,500,000 to Company A and \$2,500,000 to Company B. The
default probability is 5\% each. Assume LGD $= 100\%$. Calculate the expected loss, the
concentration ratio, and the joint probability of default.
\tcblower
\small
\textbf{Concentration ratio} $= 1/2 = \mathbf{0.5}$

\medskip
\textbf{Joint default probability and expected loss, at three correlations:}
\begin{center}
\begin{tabular}{c l r}
\arrayrulecolor{FRMRule}\toprule
\rowcolor{FRMNavy} \hdr{$\rho_{AB}$} & \hdr{$P(AB)$} & \hdr{$EL = P(AB) \times LGD \times EAD$}\\
$1.0$ & $1.0\sqrt{0.00226} + 0.0025 = 0.05$ & \$250,000\\
\rowcolor{FRMSoft}
$0.5$ & $0.5\sqrt{0.00226} + 0.0025 = 0.02627$ & \$131,350\\
$0$ & $0.05 \times 0.05 = 0.0025$ & \$12,500\\
\bottomrule
\end{tabular}
\end{center}
\end{mrexbox}

\begin{center}
\begin{tikzpicture}
\begin{axis}[
  width=10.6cm, height=4.6cm,
  xmin=-0.05, xmax=1.05, ymin=0, ymax=280000,
  xlabel={\scriptsize Default correlation $\rho_{AB}$},
  ylabel={\scriptsize Expected loss},
  xlabel style={font=\scriptsize}, ylabel style={font=\scriptsize},
  tick label style={font=\scriptsize},
  xtick={0,0.25,0.5,0.75,1},
  ytick={0,50000,100000,150000,200000,250000},
  yticklabels={\$0,\$50k,\$100k,\$150k,\$200k,\$250k},
  axis line style={draw=black!55}, clip=false]
\addplot[FRMNavy, line width=1.2pt, domain=0:1, samples=40]
  {(x*0.047539 + 0.0025)*5000000};
\node[circle, fill=FRMRed, inner sep=1.7pt] at (axis cs:0,12500) {};
\node[circle, fill=FRMRed, inner sep=1.7pt] at (axis cs:0.5,131350) {};
\node[circle, fill=FRMRed, inner sep=1.7pt] at (axis cs:1,250000) {};
\node[anchor=west, font=\scriptsize\sffamily, text=FRMRed] at (axis cs:0.03,28000)
  {\$12,500};
\node[anchor=south, font=\scriptsize\sffamily, text=FRMRed, fill=white, inner sep=1.5pt] at (axis cs:0.5,140000)
  {\$131,350};
\node[anchor=east, font=\scriptsize\sffamily, text=FRMRed] at (axis cs:0.97,232000)
  {\$250,000};
\end{axis}
\end{tikzpicture}
\end{center}
\figcap{Expected loss on the same two \$2.5m loans, plotted across the whole correlation range.
It is linear in $\rho_{AB}$ and runs from \$12,500 at zero correlation to \$250,000 at perfect
correlation --- a twentyfold difference from a single input the bank does not directly observe.}


% ==========================================================
